% Options for packages loaded elsewhere
\PassOptionsToPackage{unicode}{hyperref}
\PassOptionsToPackage{hyphens}{url}
\documentclass[
  ukrainian,
]{article}
\usepackage{xcolor}
\usepackage[margin=2.5cm]{geometry}
\usepackage{amsmath,amssymb}
\setcounter{secnumdepth}{-\maxdimen} % remove section numbering
\usepackage{iftex}
\ifPDFTeX
  \usepackage[T1]{fontenc}
  \usepackage[utf8]{inputenc}
  \usepackage{textcomp} % provide euro and other symbols
\else % if luatex or xetex
  \usepackage{unicode-math} % this also loads fontspec
  \defaultfontfeatures{Scale=MatchLowercase}
  \defaultfontfeatures[\rmfamily]{Ligatures=TeX,Scale=1}
\fi
\usepackage{lmodern}
\ifPDFTeX\else
\setmainfont{DejaVu Serif}
\setsansfont{DejaVu Sans}
\setmonofont{DejaVu Sans Mono}
\setmathfont{STIX Two Math}
\fi
\ifPDFTeX\else
  % xetex/luatex font selection
\fi
% Use upquote if available, for straight quotes in verbatim environments
\IfFileExists{upquote.sty}{\usepackage{upquote}}{}
\IfFileExists{microtype.sty}{% use microtype if available
  \usepackage[]{microtype}
  \UseMicrotypeSet[protrusion]{basicmath} % disable protrusion for tt fonts
}{}
\makeatletter
\@ifundefined{KOMAClassName}{% if non-KOMA class
  \IfFileExists{parskip.sty}{%
    \usepackage{parskip}
  }{% else
    \setlength{\parindent}{0pt}
    \setlength{\parskip}{6pt plus 2pt minus 1pt}}
}{% if KOMA class
  \KOMAoptions{parskip=half}}
\makeatother
\usepackage{longtable,booktabs,array}
\newcounter{none} % for unnumbered tables
\usepackage{calc} % for calculating minipage widths
% Correct order of tables after \paragraph or \subparagraph
\usepackage{etoolbox}
\makeatletter
\patchcmd\longtable{\par}{\if@noskipsec\mbox{}\fi\par}{}{}
\makeatother
% Allow footnotes in longtable head/foot
\IfFileExists{footnotehyper.sty}{\usepackage{footnotehyper}}{\usepackage{footnote}}
\makesavenoteenv{longtable}
\ifLuaTeX
\usepackage[bidi=basic,shorthands=off,provide=*]{babel}
\else
\usepackage[bidi=default,shorthands=off,provide=*]{babel}
\fi
\ifLuaTeX
  \usepackage{selnolig} % disable illegal ligatures
\fi
\setlength{\emergencystretch}{3em} % prevent overfull lines
\providecommand{\tightlist}{%
  \setlength{\itemsep}{0pt}\setlength{\parskip}{0pt}}
\usepackage{bookmark}
\IfFileExists{xurl.sty}{\usepackage{xurl}}{} % add URL line breaks if available
\urlstyle{same}
\hypersetup{
  pdflang={uk},
  hidelinks,
  pdfcreator={LaTeX via pandoc}}

\author{}
\date{}

\begin{document}

\section{Как играть в «Блуждающего
пьяницу»?}\label{ux43aux430ux43a-ux438ux433ux440ux430ux442ux44c-ux432-ux431ux43bux443ux436ux434ux430ux44eux449ux435ux433ux43e-ux43fux44cux44fux43dux438ux446ux443}

КРОК 1. Блукання п'яниці починаються з довільної точки сітки (у нашому
випадку --- точки А).

\begin{itemize}
\tightlist
\item
  Початкова точка - - внутрішні точки сітки - □ - Кінцева точка p;
\item
  \(g_i\) винагороди в m. \(p_I\) Фіг. 43.1. Ігрове поле для гри
  «Блукаючий п'яниця».
\end{itemize}

КРОК 2. У кожному ході гри п'яний випадково переходить до однієї з
чотирьох суміжних точок сітки (у нашому випадку сусідні точки для A ---
це точки B, C, D і E, і ймовірність досягти кожної з цих точок становить
0,25).

Таблица 43.1

{\def\LTcaptype{none} % do not increment counter
\begin{longtable}[]{@{}
  >{\raggedright\arraybackslash}p{(\linewidth - 8\tabcolsep) * \real{0.2523}}
  >{\raggedright\arraybackslash}p{(\linewidth - 8\tabcolsep) * \real{0.4541}}
  >{\raggedright\arraybackslash}p{(\linewidth - 8\tabcolsep) * \real{0.2752}}
  >{\raggedright\arraybackslash}p{(\linewidth - 8\tabcolsep) * \real{0.0092}}
  >{\raggedright\arraybackslash}p{(\linewidth - 8\tabcolsep) * \real{0.0092}}@{}}
\toprule\noalign{}
\begin{minipage}[b]{\linewidth}\raggedright
Вероятность окончания случайного блуждания в граничной точке \(p_i\) и
величина вознаграждения \(g_i\)
\end{minipage} & \begin{minipage}[b]{\linewidth}\raggedright
Граничная точка р і
\end{minipage} & \begin{minipage}[b]{\linewidth}\raggedright
\(P_{\mathcal{A}}\left(p_{i}\right)\) --- относительная доля блужданий,
закончившихся в точке \(p_{i}\)
\end{minipage} & \begin{minipage}[b]{\linewidth}\raggedright
\(g_i\) --- величина вознаграждения при достижении точки \(p_i\)
\end{minipage} & \begin{minipage}[b]{\linewidth}\raggedright
\end{minipage} \\
\midrule\noalign{}
\endhead
\bottomrule\noalign{}
\endlastfoot
1234567891011 & 0,040,150,030,060,170,050,060,150,030,060,160,04 &
110000000 & & \\
\end{longtable}
}

ti

КРОК 3. Після переходу до наступної точки процес блукання відновлюється.
Таким чином, п'яниця ходить від точки до точки, поки випадково не
опиняється на межі \(p_t\) . Тут це зупиняється, і ми записуємо номер
точки \(p_t\) . Це кінець однієї невимушеної прогулянки.

КРОК 4. Повторимо 1-3 рази достатньо і знайдемо відносну кількість
прибуттів п'яниця в кожній з меж \(p_i\) . Таблиця 1. 43.1 показує
типові результати, отримані внаслідок 100 випадкових прогулянок.

КРОК B: Припустимо, що п'яниця отримує винагороду
\(g_i = g\left(p_i\right)\) (значення граничної умови в точці \(p_i\) ,
якщо він досягає точки \(p_i\) ) і припустимо, що мета гри --- обчислити
середній \(R\left(A\right)\) винагороди за всі випадкові прогулянки,
починаючи з точки A. Цей середній виграш визначається формулою

\begin{verbatim}
Игра завершается, когда величина  R\left(A\right)  найдена. В соответствии с данными табл. 43.1 для  R\left(A\right)  получаем
\end{verbatim}

Теперь объясним, зачем нам понадобилось играть в такую игру.

\section{Для чего играть в «Блуждающего
пьяницу»?}\label{ux434ux43bux44f-ux447ux435ux433ux43e-ux438ux433ux440ux430ux442ux44c-ux432-ux431ux43bux443ux436ux434ux430ux44eux449ux435ux433ux43e-ux43fux44cux44fux43dux438ux446ux443}

Виявляється, що середня винагорода, про яку щойно йдеться, є приблизним
розв'язком проблеми Діріхле в точці A. Це цікаве спостереження базується
на двох фактах.

\begin{enumerate}
\def\labelenumi{\arabic{enumi}.}
\item
  Припустимо, що п'яниця почав свій шлях з точки, що лежить на Єрениці.
  Кожна така прогулянка одразу закінчується в одному й тому ж місці, і
  птах отримує винагороду \(g_i\) . Отже, середня винагорода за кожну
  граничну точку дорівнює \(g_i\) .
\item
  Тепер припустимо, що прогулянка починається з внутрішньої точки. Тоді
  стає зрозуміло, що середня винагорода за бал \(R(A)\) буде
  арифметичним середнім середніх нагород за чотири сусідні бали
\end{enumerate}

\[
\begin{cases} u(x, 0) = \sin(\pi x), & 0 \le x \le 1, \\ u_t(x, 0) = 0. & 0 \le x \le 1, \end{cases}
\]

Давайте знову запитаємо себе: чому середнє значення винагороди \(R(A)\)
наближається до розв'язання задачі Діріхле в точці A? Ми бачили, що
значення \(R(A)\) задовольняє двох

уравнениям

(во внутренних точках), \(R(A)\) = \(g_i\) (в граничных точках).

Если \(g_i\) ---это значения функции \(g(x, y)\) из граничного условия в
граничных точках \(p_i\) , то два наших уравнения точно совпадают с
двумя уравнениями, которые были получены при решении задачи Дирихле
методом конечных разностей. То есть величина \(R(A)\) соответствует
величине \(u_{i,f}\) в разностных уравнениях

(i, j)---внутренняя точка, \(u_{ij} = g_{ij},\) \(g_{ij}\) ---значение
решения в граничной точке (i, j).

Отже, значення \(R(A)\) дійсно наближає розв'язок рівняння часткових у
точці A.

Тепер можна вважати, що гра «Блукаючий п'яниця» складається з трьох
кроків.

\section{Розв'язання рівняння Лапласа за методом
Монте-Карло}\label{ux440ux43eux437ux432ux44fux437ux430ux43dux43dux44f-ux440ux456ux432ux43dux44fux43dux43dux44f-ux43bux430ux43fux43bux430ux441ux430-ux437ux430-ux43cux435ux442ux43eux434ux43eux43c-ux43cux43eux43dux442ux435-ux43aux430ux440ux43bux43e}

\begin{itemize}
\tightlist
\item
  Наступні правила дозволяють отримати розв'язок в одній внутрішній
  точці квадрата.
\item
  КРОК 1. Давайте зробимо кілька випадкових прогулянок, починаючи з
  точки A і закінчуючи однією з прикордонних точок. Давайте
  відстежувати, скільки разів прогулянки закінчувалися на кожній межі.
\item
  КРОК 2. Після завершення всіх прогулянок для кожної граничної точки ми
  обчислюємо відносну кількість прогулянок, які закінчилися в цій точці.
  Позначимо ці значення \(P_A(p_i)\) .

  \begin{itemize}
  \tightlist
  \item
    КРОК 3. Обчислюємо наближений розв'язок \(u(A)\) за формулою
    \(u(A) = g_1 P_A(p_1) + g_2 P_A(p_2) + \dots + g_N P_A(p_N)\) ,
  \end{itemize}
\end{itemize}

де \(g_i\) --- значення функції g у точці \(p_i\) , а N --- кількість
граничних точок.

Блукаючий п'яниця можна покращити, щоб розв'язати складнішу задачу.
Нижче наведено приклад такого оновлення.

Розв'язок задачі Діріхле з змінними коефіцієнтами.

Рассмотрим эллиптическую краевую вадачу в квадрате

Дайте физическую интерпретацию этой задачи.

\begin{enumerate}
\def\labelenumi{\arabic{enumi}.}
\setcounter{enumi}{3}
\tightlist
\item
  Для уравнения \(u_{tt} = u_{xx}\)найдите все решения вида
\end{enumerate}

\(u_{xx} + (\sin x) u_{yy} = 0\) , \(0 < x < \pi\) , \(0 < y < \pi\) ,
\((\Gamma Y)\) \(u(x, y) = g(x, y)\) на границе квадрата.

Щоб розв'язати цю задачу, замінимо \(u_{xx}\) , \(u_{yy}\) та \(\sin x\)
наступним чином:

\$\$

Будет ли решением сумма двух решений такого вида?

\end{document}
